\chapter{Introduction}

\section{Background and motivation}
\quad LEO satellite networks have gained significant attention in recent years due to their potential to provide global broadband coverage with low latency. Operating at altitudes between 160 km and 2,000 km above the Earth's surface, LEO satellites offer several advantages over GEO satellites, including reduced signal propagation delay and lower power requirements for communication \cite{lei2019broadband}. LEO satellite networks are also critical for military communications, where reliable and low-latency connectivity is essential for tactical operations, such as real-time battlefield coordination control \cite{military}. 

LEO satellites are categorized into transparent types, connecting terrestrial gNBs and GWs to UEs, and regenerative types with onboard gNBs. However, the high orbital velocity of both LEO satellites, approximately 7--8 km/s, introduces significant challenges in maintaining stable communication links such as feeder link and service link. For almost transparent types, LEO satellite feeder link operating with a break-before-make policy must perform switchovers as they pass over multiple ground GWs, which in turn induces HO for the UEs receiving LEO satellite services because the gNB associated with the each GW \cite{3gppfeeder}. In the 6G mMTC scenarios, HO failures may occur due to the transmission of HO commands to a large number of UEs. These frequent FLSO and UE HO can lead to service quality degradation and increased latency, posing a major challenge for NTN in the context of 6G communications \cite{ho_latency}. According to \cite{ho_battery}, UE power consumption during HOs ranges from 100 mW in the best case to 400 mW in the worst case. This energy overhead significantly impacts low-power battlefield devices, including IoT sensors, UAVs, and robots, hastening battery depletion and degrading network energy efficiency. Frequent FLSO and HO can disrupt mission-critical communications, increasing latency and risking security breaches in contested environments. Consequently, optimizing the FLSO strategy to minimize UE HO is crucial for ensuring seamless service in NTN.

\section{Contributions}
\quad In this paper, we propose a reinforcement learning-based FLSO scheme that optimizes HO for UEs in LEO satellite coverage. The main contributions are summarized as follows:
\begin{itemize}
    \item I created a dynamic system environment to facilitate the reinforcement learning agent's training process. By leveraging locally observable metrics, such as received signal strength and LEO satellite ephemeris data, the agent adapts to rapid orbital dynamics, thereby enhancing link stability and reducing service outages.
    \item I designed a reward function leveraging RSRP and ToS to optimize long-term LEO satellite and UE connectivity and optimize FLSO, UFLSO and UE HO in LEO satellite networks.
    \item I proposed a dueling DQN-based FLSO algorithm to enhance robustness in noisy and dynamic LEO environments for transparent LEO satellites. This algorithm dynamically selects ground GW to reduce UE HO frequency, enhance RSRP, and extend ToS across varying UE distributions.
\end{itemize}

\section{Overview}
\quad The remainder of this thesis is structured as follows: 
Sec. \uppercase\expandafter{\romannumeral2} discusses handover, feeder link switching, and related research in beam hopping and resource allocation.
Sec. \uppercase\expandafter{\romannumeral3} presents the proposed reinforcement learning framework in detail. 
Sec. \uppercase\expandafter{\romannumeral4} discusses the simulation results and compares the proposed strategy with conventional approaches. Finally, 
Sec. \uppercase\expandafter{\romannumeral5} concludes the thesis.
